\documentclass[12pt]{article}

\usepackage{sbc-template}
\usepackage{graphicx,url}
\usepackage[utf8]{inputenc}
% citações autor-ano
\usepackage[authoryear]{natbib}

\sloppy

\title{Controle de um Robô com Linguagem Natural usando LBML e Aprendizado de Máquina}

\author{Guilherme Rosa\inst{1} }

\address{Programa de Pós-Graduação em Computação Aplicada -- UDESC\\
  Joinville, SC -- Brasil
  \email{guilherme.mr@udesc.br}
}

\begin{document}

\maketitle

\begin{abstract}
Este trabalho apresenta um sistema para controlar o robô LBot por meio de comandos em linguagem natural. A solução utiliza a linguagem de domínio específico LBML e um modelo de aprendizado de máquina para traduzir comandos em português.
\end{abstract}

\begin{resumo}
O artigo descreve um sistema que permite controlar o robô LBot usando comandos em linguagem natural. A abordagem emprega a linguagem LBML e um modelo de aprendizado de máquina para traduzir comandos em português, com avaliação de desempenho por métricas de tradução automática.
\end{resumo}

\section{Introdução}
Atualmente, dois campos que despertam grande interesse tanto na academia quanto na indústria são a Robótica e a Inteligência Artificial. O avanço nessas áreas tem proporcionado impactos significativos na sociedade, permitindo a criação de soluções inovadoras que já fazem parte do cotidiano.

O projeto apresentado neste artigo integra os campos de Robótica e Inteligência Artificial. Para isso, foi desenvolvido o robô LBot, cuja principal funcionalidade é a movimentação no espaço físico. A fim de possibilitar que esses movimentos sejam descritos de forma padronizada, foi criada uma linguagem de domínio específico (DSL), denominada LBot Movement Language (LBML).

Entretanto, um desafio recorrente em aplicações de robótica é a necessidade de que o usuário aprenda comandos técnicos ou linguagens específicas para interagir com o robô. Para contornar esse problema, este trabalho propõe o uso de um modelo de aprendizado de máquina capaz de traduzir comandos em português para a LBML, permitindo que usuários interajam com o LBot por meio de linguagem natural.

Diante desse contexto, este artigo apresenta a criação da LBML, o desenvolvimento do conjunto de dados de treinamento e a implementação de um modelo baseado em arquiteturas de transformers utilizando a biblioteca PyTorch, permitindo a tradução de comandos em português para a LBML. A combinação de linguagens de domínio e aprendizado de máquina demonstra como a interação entre humanos e robôs pode ser simplificada. Além disso, a utilização de comandos em linguagem natural torna o LBot mais acessível a usuários sem conhecimento técnico, potencializando aplicações práticas em educação, entretenimento e pesquisa em robótica.

\section{Revisão de Conceitos e Trabalhos Relacionados}

O uso de \textbf{Linguagens de Domínio Específico (DSLs)} tem se consolidado como uma abordagem eficiente para simplificar a interação entre humanos e sistemas complexos, oferecendo uma alternativa à programação tradicional que exige conhecimento técnico aprofundado. DSLs são projetadas para resolver problemas de forma específica dentro de um domínio particular, possibilitando que usuários expressem suas intenções de maneira intuitiva e reduzindo a ocorrência de erros de implementação. A criação de DSLs envolve geralmente a definição de \textbf{sintaxe, semântica e regras de mapeamento} para uma linguagem formal ou plataforma de execução, permitindo que instruções de alto nível sejam traduzidas automaticamente em comandos executáveis.

No contexto de robótica, é desejável que comandos de linguagem natural possam ser convertidos para uma DSL que represente \textbf{movimentos e ações do robô} de forma precisa. A linguagem \textbf{LBot Movement Language (LBML)}, desenvolvida para o robô LBot, atende exatamente a essa necessidade, servindo como camada intermediária entre a linguagem natural do usuário e os comandos executáveis pelo sistema robótico. LBML é uma \textbf{DSL textual} que descreve movimentos no espaço físico, como avançar, girar ou manipular objetos, de maneira padronizada, garantindo que instruções complexas possam ser executadas de forma correta pelo robô.

A aplicação de DSLs em diferentes domínios já foi extensivamente estudada na literatura. \citet{ouaddi2024} propuseram uma DSL gráfica para acelerar o desenvolvimento de \emph{chatbots}, permitindo que elementos como \texttt{Intents}, \texttt{Rules}, \texttt{Stories} e \texttt{Actions} fossem definidos por meio de uma interface intuitiva de \emph{drag-and-drop}. Essa abordagem demonstrou simplificação do processo de design e codificação, reduzindo erros e custos associados à dependência de APIs externas de NLP. De forma complementar, \citet{widyani2024} exploraram o uso do modelo \textbf{GPT-3.5} para gerar DSLs de gerenciamento de dados em plataformas \emph{low-code}, mostrando que uma DSL intermediária reduziu em cerca de 77\% o uso de \texttt{tokens} em operações complexas, além de permitir tradução eficiente de instruções de linguagem natural para comandos executáveis. \citet{wu2025} aplicaram DSLs como camada intermediária na formulação de estratégias financeiras, utilizando \textbf{In-Context Learning (ICL)} para converter instruções de linguagem natural em comandos executáveis em linguagens de programação de propósito geral, destacando a importância de DSLs para aumentar interpretabilidade e reduzir redundância ou imprecisão em domínios sensíveis.

Um paralelo relevante pode ser traçado com a área de \textbf{tradução de linguagem natural para SQL (NL-to-SQL)}, que busca permitir que usuários não técnicos interajam com bancos de dados complexos. Estudos como \citet{sharma2024} compararam diferentes LLMs de código aberto em benchmarks como \textbf{SPIDER, WikiSQL e BIRD}, mostrando que modelos bem ajustados podem gerar consultas corretas mesmo em esquemas complexos e dados ruidosos. \citet{banerji2025} adaptaram LLMs para consultas em ambientes IoT industriais, utilizando mapeamento baseado em \emph{tags}, obtendo precisão e recall superiores a 0,87. De forma semelhante, \citet{chinta2025} implementaram um sistema de IA generativa para gerar consultas SQL executáveis em nuvem, alcançando até 96\% de precisão em \emph{benchmarks}, reforçando a eficácia de DSLs como intermediárias em tarefas de tradução de linguagem natural.

Além de facilitar o mapeamento de linguagem natural para comandos formais, os trabalhos relacionados também demonstram o papel de \textbf{modelos de aprendizado de máquina} na automação desse processo. O uso de \textbf{Transformers}, em especial arquiteturas \textbf{seq2seq}, tem se mostrado eficaz na tradução de sequências de texto, oferecendo capacidade de aprendizado de padrões complexos de linguagem e generalização para novos comandos. Essa abordagem permite que um modelo treinado, mesmo com um \textbf{conjunto de dados relativamente pequeno}, seja capaz de mapear instruções em linguagem natural para a representação formal na DSL, preservando o significado semântico e a correção sintática. Técnicas como \textbf{tokenização especializada} e \textbf{uso de tokens reservados} para elementos da DSL ajudam a aumentar a precisão da tradução, reduzindo ambiguidades comuns em tarefas de NL $\rightarrow$ DSL.


Apesar dos avanços em NL-to-SQL e em DSLs para chatbots, finanças e IoT, \textbf{a aplicação direta em robótica ainda é limitada}. A tradução de comandos de linguagem natural para movimentos robóticos exige precisão temporal e espacial, o que torna a construção de um modelo confiável mais desafiadora. Nesse contexto, o desenvolvimento do \textbf{LBML} e do modelo baseado em \textbf{Transformers seq2seq} apresentado neste trabalho representa uma contribuição significativa: ao criar uma linguagem intermediária padronizada e um modelo capaz de traduzir português para comandos executáveis pelo robô, este estudo avança na simplificação da interação humano-robô e amplia a acessibilidade de sistemas robóticos para usuários sem conhecimento técnico, possibilitando aplicações práticas em \textbf{educação, entretenimento e pesquisa em robótica}.
% Bibliografia
\bibliographystyle{sbc}
\bibliography{sbc-template}

\end{document}